\section{\texorpdfstring{第~\ref{chap:relativity}~章}{相对论}}\label{sec:solutions-relativity}
\subsection{\texorpdfstring{第~\ref{sec:galileo-spacetime}~节}{Galileo 时空}}\label{subsec:solutions-galileo-spacetime}

\begin{description}[font=\normalfont,leftmargin=2em,labelwidth=1.5em]
\RGsolution{ex:lucas-transatlantic-ships}{lucas-transatlantic-ships}
将横渡大西洋的轮船的运动表示为二维 Galileo 时空中的曲线, 解答就变得很简单 (参见 Fig.~\ref{fig:lucas-transatlantic-ships-spacetime}). 因此, 每艘横渡大西洋的轮船在海上会遇到另外 13 艘轮船, 相遇发生在航行中每天的午夜和正午. 若允许轮船在到达港停留一天, 用于卸货, 加油和重新装货, 就能用 15 艘横渡大西洋的轮船维持这条航线.
\RGfigure[.37\textwidth]{lucas-transatlantic-ships-spacetime}{Lucas 问题的时空图}

\RGsolution{ex:straight-line-preserving-bijection-affine}{straight-line-preserving-bijection-affine}
复合一个适当的平移后, 不妨设 $f$ 把原点映到原点. 所以只需证明 $f$ 是线性的.

首先注意, $f$ 将二维平面双射到二维平面. 实际上, 任取二维平面 $\Pi\subset\mathbb{R}^n$, 并考虑 $\Pi$ 上两两相交的 3 条直线. 它们的像也必定是两两相交的直线, 因而确定一个二维平面 $\Pi'$. $\Pi$ 中任一直线都与那 3 条两两相交的直线中的至少 2 条相交, 故其像亦然, 因而包含在 $\Pi'$ 中. 所以 $f(\Pi)\subset\Pi'$. 同样的论证给出 $f^{-1}(\Pi')\subset\Pi$, 因此 $f$ 将 $\Pi$ 满射到 $\Pi'$. 因为 $f$ 必是单射, 所以它是双射.

考虑 $f$ 在二维平面 $\Pi$ 上的限制. 由于它是双射, 它必将平行直线映成平行直线, 所以将平行四边形映为平行四边形. 因为 $f$ 把原点映到原点, 可知对任意两个向量 $v,w\in\mathbb{R}^n$, 有 $f(v+w)=f(v)+f(w)$.

最后, 考虑 $\Pi$ 上两条平行直线, 并画出与二者等距的平行直线. 任何有两边位于最初两条直线上的平行四边形, 其对角线都相交于第三条直线上的一点. 因为 $f$ 为双射且保持直线的平行性, 其像也具有这一性质. 所以 $f$ 把等距的平行直线映为等距的平行直线. 容易利用这一事实证明, $f$ 在每个二维平面上连续. 因为它可加, 所以它必是线性的.
\end{description}
